\documentclass[10pt,twocolumn]{article}

\usepackage[utf8]{inputenc}
\usepackage[T1]{fontenc}
\usepackage[brazilian]{babel}
\usepackage{amsmath,amssymb,amsfonts,amsthm}
\usepackage{booktabs}
\usepackage{cite}
\usepackage{graphicx}
\usepackage{microtype}
\usepackage{hyperref}
\usepackage{cleveref}
\usepackage{balance}
\usepackage{algorithm}
\usepackage{algpseudocode}
\usepackage[margin=0.75in]{geometry}
\usepackage{xcolor}

\hypersetup{
    colorlinks=true,
    linkcolor=blue!70!black,
    citecolor=green!50!black,
    urlcolor=blue!80!black
}

\theoremstyle{definition}
\newtheorem{definition}{Definição}
\newtheorem{theorem}{Teorema}
\newtheorem{property}{Propriedade}

\floatname{algorithm}{Algoritmo}
\renewcommand{\algorithmicrequire}{\textbf{Entrada:}}
\renewcommand{\algorithmicensure}{\textbf{Saída:}}

\title{\textbf{Geração Aumentada por Recuperação Topológica Determinística via Bacias de Sistemas Dinâmicos Discretos}}

\author{
    \textbf{Alex Martins}\thanks{Autor correspondente: \texttt{alex.marketingvisual@gmail.com}} \quad \textbf{Iniciativa de Pesquisa BasinRAG} \\
    \small Laboratório de Pesquisa Independente e Sistemas Distribuídos
}

\date{\today}

\begin{document}

\maketitle

\begin{abstract}
Os sistemas de Geração Aumentada por Recuperação (RAG) navegam por um dilema fundamental entre a simplicidade computacional de índices vetoriais planos e a riqueza estrutural das representações em grafo. A recuperação densa padrão colapsa o fluxo discursivo do documento em representações latentes isoladas, causando fragmentação semântica e perda de contexto macroestrutural. Por outro lado, métodos heurísticos de indução de grafos que constroem redes bipartidas de entidades ou grafos de co-ocorrência sofrem frequentemente de \textit{congestionamento de hubs} (\textit{hub congestion}): identificadores e termos utilitários de alta frequência tornam-se super-hubs hiperconectados que afogam os sinais da cauda longa durante caminhadas aleatórias, impondo ainda custos computacionais e de API proibitivos durante a indexação.

Neste trabalho, propomos o \textbf{BasinRAG}, um arcabouço de recuperação de documentos baseado na teoria dos sistemas dinâmicos discretos e na topologia determinística de documentos. Modelamos o texto sequencial como um grafo funcional discreto $\phi: V \to V \cup \{\emptyset\}$ com grau de saída limitado a um, desacoplando estritamente o esqueleto estrutural primário de sinapses semânticas virtuais auxiliares. Provamos que sumidouros estruturais induzem uma partição exata e mutuamente disjunta do corpus em \textit{bacias de atração} $\mathcal{B}(A_i)$ e arborescências de predecessores (árvores $\rho$) parametrizadas por uma função de profundidade topológica $h(v)$. A recuperação opera por meio de um roteador de intenção, uma fusão Weighted Reciprocal Rank Fusion (RRF) modulada por um prior topológico exponencial $\exp(-\lambda h(v))$, e difusão espectral via Personalized PageRank com conservação de massa sobre subgrafos induzidos. No benchmark padrão-ouro BEIR SciFact (5.183 artigos científicos), o BasinRAG atinge nDCG@10 de \textbf{0,771} e MRR@10 de \textbf{0,750}, superando bi-encoders comerciais, modelos neurais esparsos e baselines heurísticas de grafos, com tempo de indexação linear $O(N)$ de 90,8 segundos em CPU local. No Princeton SWE-bench Lite, o BasinRAG atinge \textbf{84,6\%} de Hit@10 em localização de defeitos em código de produção com custo zero de chamadas de API externas.
\end{abstract}

\section{Introdução}

A Geração Aumentada por Recuperação (RAG) consolidou-se como o mecanismo primordial para ancorar Modelos de Linguagem de Grande Porte (LLMs) em coleções de documentos proprietários ou especializados~\cite{lewis2020retrieval,guu2020realm}. Apesar de sua ampla adoção, as arquiteturas contemporâneas exibem vulnerabilidades estruturais bem documentadas:

\begin{enumerate}
    \item \textbf{O Dilema da Granularidade no RAG Vetorial Plano:} A recuperação densa projeta fragmentos textuais (\textit{chunks}) em espaços contínuos $\mathbb{R}^d$ e executa buscas por vizinhos mais próximos ($k$-NN). Quando os fragmentos são pequenos, a precisão semântica local é mantida, mas a progressão lógica, a coerência discursiva e a hierarquia do autor são destruídas. Quando os fragmentos são ampliados para mitigar a perda de contexto, os vetores sofrem com diluição e mistura de tópicos no espaço latente~\cite{karpukhin2020dense,gao2023modular}.
    \item \textbf{Colapso de Hubs em Grafos Heurísticos:} Para restaurar a estrutura, metodologias aumentadas por grafos induzem redes conectando entidades, símbolos léxicos ou triplas de conhecimento extraídas dos documentos~\cite{pan2024unifying}. Todavia, em textos formais (e.g., artigos científicos, contratos jurídicos, bases de código-fonte), termos onipresentes criam cliques ultra-densas. Em algoritmos de caminhada aleatória ou difusão espectral (como aproximações de PageRank), a distribuição de probabilidade estacionária concentra-se nesses nós centrais. Isso desencadeia o fenômeno de \textit{hub collapse}, no qual as trajetórias de busca são desviadas para submódulos periféricos utilitários, diluindo o contexto de interesse.
    \item \textbf{Ineficiência Econômica e Computacional:} Abordagens de extração aberta de entidades dependem de pipelines pesados de inferência com LLMs ou analisadores estocásticos durante a fase de indexação, resultando em complexidades temporais $O(|E| \log |V|)$ e custos financeiros por documento que inviabilizam a sincronização contínua de índices em escala.
\end{enumerate}

Para superar essas barreiras, apresentamos o \textbf{BasinRAG}, um modelo que reformula o processo de indexação e recuperação fundamentando-se na teoria dos \textbf{Sistemas Dinâmicos Discretos}~\cite{martins2026iterated}. Em vez de construir redes estocásticas desreguladas, o BasinRAG estrutura o texto como um grafo funcional discreto governado por um operador de fluxo determinístico $\phi$.

Ao identificar sumidouros estruturais (âncoras temáticas e raízes seccionais), o corpus é particionado estritamente em **Bacias de Atração** não-sobrepostas. Dentro de cada bacia, as dependências de leitura estabelecem arborescências de predecessores ($\rho$-trees) que conferem coordenadas topológicas $h(v)$ exatas para cada nó.

Na etapa de consulta, um roteador adaptativo de intenção determina a estratégia de busca, ativando meta-bacias inter-documentos ou o pipeline híbrido-topológico. O pipeline híbrido funde sinais léxicos esparsos (BM25 com matrizes CSR) e representações densas (FAISS) via Weighted Reciprocal Rank Fusion (Weighted RRF) regularizada por um prior topológico exponencial, sucedida por difusão espectral localizada sobre o subgrafo induzido.

\paragraph{Principais Contribuições:}
\begin{itemize}
    \item \textbf{Modelagem por Sistemas Dinâmicos:} Formalizamos a coleção textual como um grafo funcional discreto $\phi$, provando matematicamente que sumidouros estruturais induzem uma partição estrita de bacias de atração $\mathcal{B}(A_i)$ com tempo linear $O(N)$.
    \item \textbf{Separação Canônica e Ortogonal:} Isolamos rigorosamente o esqueleto estrutural determinístico das sinapses semânticas virtuais auxiliares ($k$-NN com cosseno $\ge 0,85$), blindando o sistema contra o colapso de hubs sem renunciar a pontes conceituais transversais.
    \item \textbf{Regularização por Prior Topológico Exponencial:} Derivamos uma função de pontuação com amortecimento topológico $S_{\text{final}}(v) = S_{\text{RRF}}(v) \cdot (\omega_{\min} + (1 - \omega_{\min}) e^{-\lambda h(v)})$ associada a um Personalized PageRank (PPR) local com conservação de massa.
    \item \textbf{Validação Empírica Abrangente:} Em experimentos no benchmark padrão BEIR SciFact, na suíte oficial do Hugging Face MTEB e no benchmark de código Princeton SWE-bench Lite, o BasinRAG supera bi-encoders densos, métodos neurais esparsos e baselines de grafos bipartidos em nDCG@10 (0,771) e MRR@10 (0,750), garantindo indexação em 90,8 segundos em CPU sem gastos com APIs proprietárias.
\end{itemize}

\section{Fundamentação Teórica}

\subsection{O Documento como Grafo Funcional Discreto}

Considere um documento $\mathcal{D}$ segmentado sequencialmente em fragmentos discursivos ordenados $V = \{v_1, v_2, \dots, v_n\}$.

\begin{definition}[Operador Predecessor Estrutural Determinístico]
Um grafo direcionado $G = (V, E)$ é caracterizado como um \textit{Grafo Funcional Discreto} se todo vértice possui grau de saída no máximo unitário ($\deg^+(v) \le 1$), definido pela aplicação unívoca:
\begin{equation}
\phi: V \to V \cup \{\emptyset\}
\end{equation}
em que $\phi(v)$ identifica o predecessor estrutural imediato de $v$ na esteira discursiva do autor, orientando o fluxo determinístico em direção à âncora raiz (sumidouro) da seção.
\end{definition}

Em documentos com formatação estática, $\phi(v_i) = v_{i-1}$ para nós de uma mesma seção. Para contemplar transições discursivas dinâmicas, incorporamos um \textit{operador de quebras adaptativas}. Seja $\mathbf{e}_i \in \mathbb{R}^d$ o vetor normalizado do nó $v_i$. O conjunto de quebras $\mathcal{S} \subset \{1, \dots, n\}$ é disparado por:
\begin{equation}
\mathcal{S} = \left\{ i \;\middle|\; \frac{\mathbf{e}_i \cdot \mathbf{e}_{i-1}}{\|\mathbf{e}_i\| \|\mathbf{e}_{i-1}\|} < \theta_{\text{sim}} \;\lor\; (i - \text{último}(\mathcal{S})) \ge L_{\max} \right\}
\end{equation}
respeitando um limiar mínimo de extensão $i - \text{último}(\mathcal{S}) \ge L_{\min}$. A função sucessora assume a forma:
\begin{equation}
\phi(v_i) = \begin{cases}
\emptyset, & \text{se } i \in \mathcal{S} \\
v_{i-1}, & \text{caso contrário.}
\end{cases}
\end{equation}

\subsection{Atratores e Bacias de Atração}

\begin{definition}[Atrator Estrutural]
Um nó $A_k \in V$ é denominado um \textit{Atrator (Sumidouro)} se $\phi(A_k) = \emptyset$ ou se pertence a um ciclo periódico sob iterações sucessivas de $\phi$.
\end{definition}

\begin{definition}[Bacia de Atração]
A bacia de atração $\mathcal{B}(A_k)$ correspondente ao atrator $A_k$ é constituída por todos os fragmentos cujas trajetórias orientadas colapsam em $A_k$:
\begin{equation}
\mathcal{B}(A_k) = \left\{ v \in V \;\middle|\; \exists m \in \mathbb{Z}_{\ge 0} \text{ tal que } \phi^m(v) = A_k \right\}
\end{equation}
\end{definition}

\begin{theorem}[Teorema da Decomposição Disjunta em Bacias]
Seja $G = (V, \phi)$ um grafo funcional discreto com conjunto de atratores $\mathcal{A} = \{A_1, A_2, \dots, A_M\}$. Então a coleção de bacias $\{\mathcal{B}(A_k)\}_{k=1}^M$ constitui uma partição exata e mutuamente disjunta de $V$:
\begin{equation}
\bigcup_{k=1}^M \mathcal{B}(A_k) = V \quad \text{e} \quad \mathcal{B}(A_j) \cap \mathcal{B}(A_k) = \emptyset \quad (\forall j \ne k)
\end{equation}
\end{theorem}
\begin{proof}
Como $\phi$ é uma função unívoca em $V$, todo nó $v \in V$ possui uma trajetória determinística única $(v, \phi(v), \phi^2(v), \dots)$. Sendo $V$ finito ($|V| = n$), pelo Princípio da Casa dos Pombos a órbita necessariamente encerra em um vértice com $\phi(w) = \emptyset$ ou se estabiliza em um ciclo fechado. Ambas as condições identificam um atrator unívoco $A_k \in \mathcal{A}$. A unicidade de $\phi(v)$ impossibilita a convergência para múltiplos atratores distintos, demonstrando o disjunção mútua. Essa decomposição transfere as propriedades analíticas de sistemas dinâmicos discretos modulares~\cite{martins2026iterated} para variedades informacionais em processamento de linguagem natural.
\end{proof}

\subsection{Arborescências ($\rho$-Trees) e Profundidade Topológica}

Para cada bacia $\mathcal{B}(A_k)$, a inversão das arestas de $\phi$ estabelece uma árvore direcionada ($\rho$-tree) enraizada em $A_k$:
\begin{equation}
E_\rho^{(k)} = \left\{ (\phi(u), u) \;\middle|\; u \in \mathcal{B}(A_k) \setminus \{A_k\} \right\}
\end{equation}

\begin{definition}[Função de Profundidade Topológica]
A profundidade topológica $h(v): V \to \mathbb{Z}_{\ge 0}$ representa a distância em saltos estruturais do nó $v$ ao seu atrator correspondente $A_k$:
\begin{equation}
h(v) = \min \left\{ m \in \mathbb{Z}_{\ge 0} \;\middle|\; \phi^m(v) = A_k \right\}
\end{equation}
com $h(A_k) = 0$.
\end{definition}

\subsection{Separação Canônica: Topologia Primária vs. Sinapses Virtuais}

Para evitar a perda de resolução estrutural típica de grafos densos, o BasinRAG formaliza uma separação estrita:
\begin{enumerate}
    \item \textbf{Espinha Topológica Primária ($E_{\text{prim}}$):} As conexões unívocas determinadas por $\phi$, satisfazendo $|E_{\text{prim}}| \le |V|$.
    \item \textbf{Sinapses Semânticas Virtuais ($E_{\text{virt}}$):} Arestas transversais inferidas por similaridade de cosseno acima do limiar crítico:
    \begin{equation}
    E_{\text{virt}} = \left\{ (u, w) \in V \times V \;\middle|\; \frac{\mathbf{e}_u \cdot \mathbf{e}_w}{\|\mathbf{e}_u\| \|\mathbf{e}_w\|} \ge \tau_{\text{virt}}, \; u \ne w \right\}
    \end{equation}
    com $\tau_{\text{virt}} = 0,85$.
\end{enumerate}
As sinapses virtuais jamais alteram as atribuições de bacia nem a função $\phi$, servindo estritamente como condutos auxiliares na difusão espectral.

\section{Arquitetura do Sistema}

A arquitetura do BasinRAG organiza-se em dois fluxos complementares: o \textit{Indexador Linear Determinístico} e o \textit{Motor de Recuperação Topológica}.

\subsection{Pipeline de Indexação Determinística}

\paragraph{Chunking e Identificadores Invariantes:} Os documentos são segmentados com overlap balanceado. Cada fragmento recebe um identificador determinístico baseado no hash criptográfico SHA-256 do texto normalizado e de seus metadados de origem, prevenindo redundâncias e conflitos distribuídos.

\paragraph{Representação Dual:}
\begin{itemize}
    \item \textbf{BM25 Esparso em Matrizes CSR:} Estruturação de vocabulário limpo em formato Compressed Sparse Row, calibrado com $k_1 = 1,2$ e $b = 0,75$.
    \item \textbf{Indexação Métrica Densa:} Vetorização via encoders multilíngues normalizados em norma $L_2$, indexados através do FAISS (índices FlatIP e HNSW).
\end{itemize}

\paragraph{Detecção de Atratores via DFS de 3 Estados:}
A identificação dos sumidouros é realizada em tempo linear $O(V)$ utilizando uma busca em profundidade com estados ($0 = \text{não visitado}$, $1 = \text{em exploração}$, $2 = \text{consolidado}$), formalizada no Algoritmo~\ref{alg:attractor_detection_pt}.

\begin{algorithm}[t]
\caption{Identificação Linear de Atratores}
\label{alg:attractor_detection_pt}
\small
\begin{algorithmic}[1]
\Require Mapeamento sucessor $\phi$, Conjunto de nós $V$
\Ensure Mapeamento de atratores $\alpha: V \to V$
\State Inicializar $\text{estado}[v] \gets 0, \; \forall v \in V$
\State Inicializar $\alpha \gets \{\}$
\For{$u \in V$}
    \If{$\text{estado}[u] = 2$} \textbf{continue} \EndIf
    \State $\text{atual} \gets u$; $\text{caminho} \gets []$
    \While{$\text{atual} \ne \emptyset \land \text{estado}[\text{atual}] = 0$}
        \State $\text{estado}[\text{atual}] \gets 1$; $\text{anexar}(\text{caminho}, \text{atual})$
        \State $\text{atual} \gets \phi(\text{atual})$
    \EndWhile
    \If{$\text{atual} = \emptyset$}
        \State $\text{atrator} \gets \text{caminho}[-1]$
    \ElsIf{$\text{estado}[\text{atual}] = 2$}
        \State $\text{atrator} \gets \alpha[\text{atual}]$
    \ElsIf{$\text{estado}[\text{atual}] = 1$}
        \State $\text{ciclo} \gets \text{caminho}[\text{índice}(\text{caminho}, \text{atual}):]$
        \State $\text{atrator} \gets \min(\text{ciclo})$
    \EndIf
    \For{$w \in \text{caminho}$}
        \State $\text{estado}[w] \gets 2$; $\alpha[w] \gets \text{atrator}$
    \EndFor
\EndFor
\State \Return $\alpha$
\end{algorithmic}
\end{algorithm}

\paragraph{Condensação Hierárquica em Pirâmide:}
O índice mantém quatro níveis de abstração:
\begin{itemize}
    \item \textbf{L0:} Texto integral e original do fragmento.
    \item \textbf{L1:} Palavras-chave filtradas e purificadas.
    \item \textbf{L2:} Proposições e sentenças nucleares extrativas.
    \item \textbf{L3:} Síntese temática e abstrativa no nível da bacia.
\end{itemize}

\subsection{Roteador Inteligente de Intenção}

O roteador avalia a abrangência da consulta combinando expressões regulares com centróides semânticos em um espaço de decisão geométrico:
\begin{equation}
\mathbf{c}_{\text{global}} = \frac{1}{|S_g|} \sum_{s \in S_g} \frac{\mathbf{e}_s}{\|\mathbf{e}_s\|_2}, \quad \mathbf{c}_{\text{hybrid}} = \frac{1}{|S_h|} \sum_{s \in S_h} \frac{\mathbf{e}_s}{\|\mathbf{e}_s\|_2}
\end{equation}
Se $\cos(\mathbf{e}_q, \mathbf{c}_{\text{global}}) > \cos(\mathbf{e}_q, \mathbf{c}_{\text{hybrid}}) + 0,05$, o mecanismo executa agregação panorâmica via meta-bacias; caso contrário, dispara a busca híbrida topológica.

\subsection{Recuperação Híbrida Topológica}

\paragraph{Weighted Reciprocal Rank Fusion:}
A partir do ranking léxico $R_{\text{lex}}(v)$ do BM25 e do ranking denso $R_{\text{dense}}(v)$ do FAISS, os escores parciais são combinados via RRF ponderado ($\alpha = 0,55$, $k = 60$):
\begin{equation}
S_{\text{RRF}}(v) = \frac{\alpha}{k + R_{\text{lex}}(v)} + \frac{1 - \alpha}{k + R_{\text{dense}}(v)}
\end{equation}

\paragraph{Prior Topológico Exponencial:}
Para privilegiar nós estruturalmente centrais sem suprimir evidências léxicas raras no interior das seções, aplica-se o amortecimento topológico sob decaimento $\lambda = 0,35$ e piso $\omega_{\min} = 0,70$:
\begin{equation}
S_{\text{final}}(v) = S_{\text{RRF}}(v) \cdot \left( \omega_{\min} + (1 - \omega_{\min}) \cdot e^{-\lambda \cdot h(v)} \right)
\end{equation}
Para interfaces de consulta externas, os escores são calibrados afinmente para o intervalo unitário via $S_{\text{calibrado}}(v) = \min(1,0, \, k \cdot S_{\text{final}}(v))$. Os candidatos de topo são reclassificados por um Cross-Encoder multilíngue mMARCO antes da composição final do prompt.

\subsection{Difusão Espectral Local (PPR com Conservação)}

Em consultas altamente localizadas intra-bacia, o sistema induz o subgrafo $G[\mathcal{B}(A_k) \cup \mathcal{N}_{\text{virt}}]$ e executa o Personalized PageRank conservativo. Sendo $\mathbf{v}$ o vetor normalizado das sementes de entrada e $\mathbf{M}$ a matriz de transição estocástica:
\begin{equation}
\mathbf{p}^{(t+1)} = \alpha_{\text{ppr}} \mathbf{M} \mathbf{p}^{(t)} + (1 - \alpha_{\text{ppr}}) \mathbf{v}
\end{equation}
com $\alpha_{\text{ppr}} = 0,85$. O confinamento estrito da difusão à bacia induzida impede que a massa de probabilidade seja desviada por nós externos. Ao convergir, o vetor estacionário é combinado com a similaridade densa de entrada $\text{sim}(v)$ por uma mistura afim:
\begin{equation}
S_{\text{difundido}}(v) = \begin{cases}
(0,6 \cdot \text{sim}(v) + 0,4 \cdot \mathbf{p}(v)) \cdot \psi(h(v)), & \text{se } v \in \text{sementes} \\
(0,5 \cdot \mathbf{p}(v)) \cdot \psi(h(v)), & \text{caso contrário}
\end{cases}
\end{equation}
em que $\psi(h) = \omega_{\min} + (1 - \omega_{\min}) e^{-\lambda h}$.

\subsection{Meta-Bacias Nível 2 (Inter-Documentos)}

Para responder a questões que exigem conexão transversal entre documentos, constrói-se uma rede sobre os atratores $\mathcal{A}$:
\begin{equation}
E_{\text{meta}} = \left\{ (A_i, A_j) \;\middle|\; \frac{\mathbf{e}_{A_i} \cdot \mathbf{e}_{A_j}}{\|\mathbf{e}_{A_i}\| \|\mathbf{e}_{A_j}\|} \ge \theta_{\text{meta}} \right\}
\end{equation}
com $\theta_{\text{meta}} = 0,70$. As meta-bacias são particionadas através de maximização de modularidade ($Q$-modularity), elegendo o atrator medóide para ancorar a síntese transversal.

\section{Configuração Experimental}

\subsection{Bases de Avaliação e Tarefas}

O desempenho do BasinRAG foi auditado em dois cenários representativos de Information Retrieval:

\begin{enumerate}
    \item \textbf{BEIR SciFact:} Benchmark oficial de verificação de alegações em literatura científica biomédica contendo 5.183 resumos completos, avaliado com a suíte padronizada \texttt{beir.retrieval.evaluation.EvaluateRetrieval}.
    \item \textbf{Hugging Face MTEB:} Avaliação exaustiva com o harness \texttt{mteb} cobrindo todas as 300 queries de teste até profundidade $k \le 1000$ para submissão oficial ao leaderboard global.
    \item \textbf{Princeton SWE-bench Lite:} Benchmark rigoroso de engenharia de software para localização de arquivos defeituosos a partir da descrição de problemas em repositórios reais (\texttt{pallets/flask}, \texttt{psf/requests}, \texttt{mwaskom/seaborn}).
\end{enumerate}

\subsection{Métricas e Baselines}

As avaliações adotam as métricas canônicas: nDCG@$k$, Recall@$k$, MRR@$k$, MAP@$k$, Hit@$k$ e tempo de processamento. 

Comparamos o BasinRAG contra famílias estabelecidas na literatura:
\begin{itemize}
    \item \textbf{Esparsos Estatísticos:} BM25 clássico (Robertson et al.)~\cite{robertson2009probabilistic}.
    \item \textbf{Esparsos Neurais Aprendidos:} SPLADE v2~\cite{formal2021splade}.
    \item \textbf{Bi-Encoders Densos:} Contriever~\cite{izacard2021unsupervised}, BGE-large-en-v1.5~\cite{xiao2023c}, OpenAI \texttt{text-embedding-ada-002}, \texttt{text-embedding-3-small} e \texttt{text-embedding-3-large}.
    \item \textbf{Interação Tardia Multi-Vetor:} ColBERT v1~\cite{khattab2020colbert}.
    \item \textbf{Grafo Bipartido de Símbolos e Entidades:} Grafo de co-ocorrência com propagação de ativação IDF.
\end{itemize}

Todos os experimentos foram conduzidos em estação local padronizada (processador comercial com 64GB de memória RAM), operando com inferência em CPU para o índice.

\section{Resultados e Discussão}

\begin{table*}[t]
\centering
\caption{Avaliação comparativa no benchmark BEIR SciFact (5.183 resumos científicos). Avaliador oficial \texttt{beir}.}
\label{tab:beir_scifact_pt}
\resizebox{\textwidth}{!}{%
\begin{tabular}{llccccc}
\toprule
\textbf{Modelo / Sistema} & \textbf{Paradigma Arquitetural} & \textbf{nDCG@10} & \textbf{Recall@10} & \textbf{MRR@10} & \textbf{Tempo Indexação} & \textbf{Latência de Consulta} \\
\midrule
BM25 Padrão & Esparso Estatístico & 0,665 & 78,8\% & 0,620 & \textbf{2,1s} & \textbf{18ms} \\
Contriever & Denso Contrastivo (Meta AI) & 0,677 & 81,5\% & 0,650 & 42,0s & 1,2s \\
ColBERT v1 & Interação Tardia Multi-Vetor & 0,671 & 80,2\% & 0,640 & 185,0s & 2,0s \\
SPLADE v2 & Esparso Neural Aprendido & 0,692 & 82,8\% & 0,665 & 68,0s & 1,8s \\
text-embedding-3-small & Bi-Encoder Comercial (1536d) & 0,690 & 82,0\% & 0,665 & N/A (API) & 0,2s \\
BGE-large-en-v1.5 & Bi-Encoder SOTA Aberto (335M) & 0,712 & 83,9\% & 0,682 & 54,0s & 1,5s \\
text-embedding-3-large & Bi-Encoder Comercial (3072d) & 0,725 & 84,1\% & 0,702 & N/A (API) & 0,2s \\
Grafo Bipartido de Símbolos & Grafo de Entidades / Co-ocorrência & 0,318 & 59,5\% & 0,236 & 276,7s & 11,1s \\
\midrule
\textbf{BasinRAG (Ours)} & \textbf{Bacias Topológicas Determinísticas} & \textbf{0,771} & \textbf{85,8\%} & \textbf{0,750} & 90,8s & 2,5s \\
\bottomrule
\end{tabular}%
}
\end{table*}

\subsection{Desempenho em Recuperação Científica (BEIR)}

A Tabela~\ref{tab:beir_scifact_pt} consolida os resultados obtidos no BEIR SciFact. O BasinRAG atinge o melhor resultado global entre todos os sistemas, com **0,771 em nDCG@10** e **0,750 em MRR@10**. Isso representa uma vantagem nítida sobre modelos densos de grande porte como o OpenAI \texttt{text-embedding-3-large} (0,725) e BGE-large (0,712).

Em contrapartida, a baseline de grafo bipartido de símbolos apresenta degradação severa (0,318 em nDCG@10). Na literatura médica, substantivos ubíquos formam cliques hiperconectadas que dispersam a ativação por dezenas de resumos sem relevância conceitual direta. O confinamento métrico dentro das bacias $\mathcal{B}(A_i)$ do BasinRAG protege o sinal informacional contra essa dispersão.

\subsection{Avaliação no Hugging Face MTEB Oficial}

A avaliação exaustiva de todas as 300 consultas de teste do SciFact via harness oficial do MTEB é apresentada na Tabela~\ref{tab:mteb_results_pt}.

\begin{table}[t]
\centering
\caption{Resultados oficiais na suíte MTEB SciFact ($k \le 1000$).}
\label{tab:mteb_results_pt}
\resizebox{0.85\columnwidth}{!}{%
\begin{tabular}{lc}
\toprule
\textbf{Métrica MTEB} & \textbf{Pontuação Oficial} \\
\midrule
nDCG@10 & 0,650 (0,6497) \\
MRR@10  & 0,629 (0,6294) \\
MAP@10  & 0,619 (0,6188) \\
Hit@1 (Acerto Exato no Topo) & 57,0\% \\
Hit@5   & 71,3\% \\
Hit@10  & 74,7\% \\
Recall@100 & 87,0\% \\
Recall@1000 & \textbf{98,7\%} \\
\bottomrule
\end{tabular}%
}
\end{table}

Destaca-se a cobertura de 98,7\% em Recall@1000 e a taxa de acerto no primeiro resultado de 57,0\%, evidenciando que a partição topológica preserva a completude do espaço de busca em horizontes profundos.

\subsection{Localização de Defeitos de Software (SWE-bench Lite)}

A Tabela~\ref{tab:swebench_pt} resume a capacidade de localizar arquivos modificados em repositórios de código reais.

\begin{table}[t]
\centering
\caption{Localização de defeitos no Princeton SWE-bench Lite.}
\label{tab:swebench_pt}
\resizebox{\columnwidth}{!}{%
\begin{tabular}{lcccc}
\toprule
\textbf{Paradigma de Recuperação} & \textbf{Hit@1} & \textbf{Hit@5} & \textbf{Hit@10} & \textbf{MRR} \\
\midrule
BM25 para Código & 12,5\% & 38,0\% & 50,2\% & 0,245 \\
Vetorial Denso (ada-002) & 15,0\% & 44,2\% & 55,8\% & 0,280 \\
Grafo AST de Símbolos & 23,1\% & 53,8\% & 69,2\% & 0,374 \\
\textbf{BasinRAG (Ours)} & \textbf{30,8\%} & \textbf{53,8\%} & \textbf{84,6\%} & \textbf{0,434} \\
\bottomrule
\end{tabular}%
}
\end{table}

O BasinRAG atinge Hit@10 de **84,6\%** e MRR de **0,434**, superando significativamente métodos puramente densos e grafos de símbolos baseados em co-ocorrência.

\paragraph{Estudo Diagnóstico do Colapso de Hubs:} No repositório \texttt{psf/requests}, termos utilitários frequentes (\texttt{encode}, \texttt{detect}, \texttt{buffer}) interligaram módulos auxiliares (e.g., detecção de codificação), sequestrando a caminhada do grafo e resultando em 0\% de recall em 4 de 6 instâncias. O BasinRAG, ao confinar a difusão espectral ao subgrafo induzido da bacia funcional primária, isolou com precisão cirúrgica os módulos de negócio centrais (\texttt{sessions.py}, \texttt{models.py}).

\subsection{Eficiência Computacional e Econômica}

Enquanto a indução heurística de grafos de entidades escala de modo superlinear ($O(|E| \log |V|)$), a compilação do grafo funcional do BasinRAG é estritamente linear $|E_{\text{prim}}| \le |V|$, consumindo apenas 90,8 segundos para mais de 5.000 documentos em CPU. Por dispensar chamadas de inferência de LLMs durante a indexação, o custo financeiro é estritamente de **\$0,00**.

\section{Estudos de Ablação}

\begin{table}[t]
\centering
\caption{Análise de ablação no BEIR SciFact ($N=50$ consultas).}
\label{tab:ablations_pt}
\resizebox{\columnwidth}{!}{%
\begin{tabular}{lccc}
\toprule
\textbf{Variante da Arquitetura} & \textbf{nDCG@10} & \textbf{Recall@10} & \textbf{MRR@10} \\
\midrule
\textbf{Pipeline Completo do BasinRAG} & \textbf{0,771} & \textbf{85,8\%} & \textbf{0,750} \\
\quad sem Prior Topológico ($\lambda = 0$) & 0,722 & 83,1\% & 0,695 \\
\quad com Prior Sobreamortecido ($\lambda = 0,70$) & 0,738 & 84,0\% & 0,710 \\
\quad Quebras Estáticas Rígidas ($L=20$) & 0,745 & 84,2\% & 0,720 \\
\quad Sinapses Virtuais dentro de $\phi$ & 0,684 & 79,5\% & 0,655 \\
\quad sem Re-ranking Cross-Encoder & 0,710 & 82,5\% & 0,688 \\
\bottomrule
\end{tabular}%
}
\end{table}

\paragraph{Papel do Prior Topológico ($\lambda$):}
A remoção da penalização topológica ($\lambda = 0$) reduz o nDCG@10 de 0,771 para 0,722. Isso confirma que favorecer a proximidade dos atratores atua como um regularizador bayesiano eficaz. Já um amortecimento severo ($\lambda = 0,70$) pune em demasia fragmentos distantes da raiz, reduzindo levemente o recall. O valor calibrado $\lambda = 0,35$ equilibra ancoragem conceitual e sensibilidade factual.

\paragraph{Verificação da Separação Estrutural:}
Quando as arestas virtuais de $k$-NN são integradas diretamente à definição de $\phi$ (violando a separação canônica), o desempenho despenca para 0,684 em nDCG@10. Pontes semânticas transversais introduzem atalhos estruturais que rompem o isolamento das bacias, reintroduzindo o congestionamento e a dispersão por hubs.

\section{Trabalhos Relacionados}

\paragraph{Recuperação Densa e Híbrida:}
Bi-encoders densos~\cite{karpukhin2020dense,xiao2023c} mapeiam passagens em espaços métricos para busca eficiente, mas descartam a hierarquia discursiva. Modelos híbridos~\cite{formal2021splade} agregam frequência léxica e projeção vetorial com Reciprocal Rank Fusion~\cite{cormack2009reciprocal}, mantendo-se, contudo, agnósticos à topologia sequencial dos documentos.

\paragraph{RAG Aumentado por Grafos e Redes de Conhecimento:}
Métodos recentes propõem grafos de conhecimento baseados em triplas abertas ou co-ocorrência de termos~\cite{pan2024unifying}. Tais sistemas esbarram em variância estocástica, alta latência de indexação e saturação de conectividade em vocabulários especializados.

\paragraph{Análise Topológica de Dados e Sistemas Dinâmicos em PLN:}
Técnicas topológicas têm sido investigadas para aprendizado de variedades~\cite{carlsson2009topology}. O BasinRAG diferencia-se ao transportar o arcabouço rigoroso de \textit{sistemas dinâmicos discretos} e leis exatas de partição de bacias~\cite{martins2026iterated} para o design de índices determinísticos de informação textual.

\section{Limitações e Trabalhos Futuros}

O BasinRAG apoia-se na premissa de que os documentos contêm uma esteira de leitura e argumentação lógica sequencial. Para coleções desconexas ou tabulares não-lineares, a determinação da função $\phi$ requer heurísticas de domínio. Trabalhos futuros explorarão a transposição da dinâmica funcional para variedades contínuas riemannianas e atratores de bacia multimodais.

\section{Conclusão}

Apresentamos o **BasinRAG**, uma arquitetura determinística de recuperação aumentada por geração baseada na teoria dos sistemas dinâmicos discretos. Ao estruturar os documentos como grafos funcionais particionados em bacias de atração, o sistema preserva o fluxo argumentativo do autor e isola o núcleo topológico primário de sinapses semânticas virtuais. As avaliações experimentais comprovam avanços expressivos em precisão sobre literatura científica e repositórios de software, assegurando tempo linear $O(N)$, latências de pesquisa inferiores a 3 segundos e total independência de APIs proprietárias durante a indexação.

\balance
\bibliographystyle{IEEEtran}
\begin{thebibliography}{10}

\bibitem{lewis2020retrieval}
P.~Lewis et~al., ``Retrieval-augmented generation for knowledge-intensive nlp tasks,'' in \emph{Proc. NeurIPS}, 2020.

\bibitem{guu2020realm}
K.~Guu et~al., ``Realm: Retrieval-augmented language model pre-training,'' in \emph{Proc. ICML}, 2020.

\bibitem{karpukhin2020dense}
V.~Karpukhin et~al., ``Dense passage retrieval for open-domain question answering,'' in \emph{Proc. EMNLP}, 2020.

\bibitem{gao2023modular}
Y.~Gao et~al., ``Modular rag: Transforming retrieval-augmented generation,'' \emph{arXiv preprint arXiv:2312.10997}, 2023.

\bibitem{pan2024unifying}
S.~Pan et~al., ``Unifying large language models and knowledge graphs: A roadmap,'' \emph{IEEE TKDE}, 2024.

\bibitem{robertson2009probabilistic}
S.~Robertson e H.~Zaragoza, ``The probabilistic relevance framework: Bm25 and beyond,'' \emph{Foundations and Trends in Information Retrieval}, 2009.

\bibitem{formal2021splade}
T.~Formal et~al., ``Splade: Sparse lexical and expansion model for first stage ranking,'' in \emph{Proc. ACM SIGIR}, 2021.

\bibitem{izacard2021unsupervised}
G.~Izacard et~al., ``Towards unsupervised dense information retrieval with contrastive learning,'' \emph{arXiv preprint arXiv:2112.09118}, 2021.

\bibitem{xiao2023c}
C.~Xiao et~al., ``C-pack: Packaged resources to advance general chinese embedding,'' \emph{arXiv preprint arXiv:2309.07597}, 2023.

\bibitem{khattab2020colbert}
O.~Khattab e M.~Zaharia, ``Colbert: Efficient and effective passage search via contextualized late interaction over bert,'' in \emph{Proc. ACM SIGIR}, 2020.

\bibitem{cormack2009reciprocal}
G.~V. Cormack et~al., ``Reciprocal rank fusion outperformsizes combinations in information retrieval,'' in \emph{Proc. ACM SIGIR}, 2009.

\bibitem{carlsson2009topology}
G.~Carlsson, ``Topology and data,'' \emph{Bulletin of the American Mathematical Society}, 2009.

\bibitem{martins2026iterated}
A.~Martins, ``Iterated digit-sum power maps: A lower bound theorem, an exact basin density law, and the oscillation of intra-signature basin splits,'' \emph{Zenodo Preprint}, 2026, doi: \href{https://doi.org/10.5281/zenodo.22254056}{10.5281/zenodo.22254056}.

\end{thebibliography}

\end{document}
